```latex
\chapter{REVUE DE LA LITTÉRATURE ET FONDEMENTS THÉORIQUES}
\label{chapitre-i-revue-de-la-litterature-et-fondements-theoriques}

\section{Concepts fondamentaux liés à la distribution de l'eau potable}
\label{concepts-fondamentaux-lies-a-la-distribution-de-leau-potable}

\subsection{Définition de l'eau potable}
\label{definition-de-leau-potable}

L'eau potable désigne une eau dont les caractéristiques physiques,
chimiques, biologiques et microbiologiques respectent les exigences
nécessaires à la consommation humaine sans risque significatif pour la
santé. Sa sécurité doit être maintenue tout au long de la chaîne
d'approvisionnement, depuis le captage jusqu'au consommateur
\cite{ref4}, \cite{ref23}.

L'accès à une eau potable sûre constitue également un droit humain
fondamental reconnu par les Nations Unies et participe à la réalisation
de l'Objectif de développement durable 6 (ODD 6), consacré à l'eau
potable et à l'assainissement \cite{ref3}. Dans cette perspective, la
qualité de l'eau dépend non seulement du traitement, mais également de
la fiabilité des infrastructures chargées de son transport et de sa
distribution \cite{ref1}.

\subsection{Réseau de distribution d'eau potable}
\label{reseau-de-distribution-deau-potable}

Le réseau de distribution d'eau potable est l'ensemble des ouvrages et
équipements permettant d'acheminer l'eau traitée jusqu'aux points de
consommation. Il comprend notamment les conduites, les réservoirs, les
stations de pompage, les vannes, les dispositifs de contrôle et les
équipements de mesure \cite{ref11}.

Le fonctionnement du réseau dépend de l'interaction entre ces différents
éléments. Une défaillance d'une conduite, d'une pompe, d'une vanne ou
d'un réservoir peut entraîner une baisse de pression, une interruption
de service ou une perte d'eau. La surveillance de paramètres tels que
la pression, le débit et le niveau des réservoirs constitue donc une
condition importante pour assurer la continuité du service \cite{ref20}.

\subsection{Composantes du réseau}
\label{composantes-du-reseau}

Les principales composantes d'un réseau d'alimentation en eau potable
sont les ouvrages de captage, les stations de traitement, les réservoirs,
les conduites, les stations de pompage et les vannes. Les conduites
assurent le transport de l'eau, les réservoirs permettent son stockage et
la régulation du système, tandis que les pompes et les vannes contribuent
au maintien des conditions hydrauliques nécessaires à la distribution
\cite{ref33}, \cite{ref34}, \cite{ref35}.

Dans les réseaux modernes, des capteurs peuvent être associés à ces
composants afin de mesurer différents paramètres de fonctionnement.
Ces informations constituent une base importante pour la surveillance
numérique et l'analyse de l'état du réseau \cite{ref20}.

\section{Déploiement d'un réseau de distribution d'eau potable}
\label{deploiement-dun-reseau-de-distribution-deau-potable}

Le déploiement d'un réseau d'eau potable comprend la planification, la
conception, le dimensionnement et l'installation des infrastructures
nécessaires à la desserte d'une population. Il doit tenir compte des
besoins en eau, de la topographie, de la répartition de la population,
des capacités des infrastructures existantes et des possibilités
d'extension du réseau \cite{ref38}, \cite{ref39}.

La planification spatiale permet de déterminer la localisation des
conduites, des réservoirs et des équipements hydrauliques. Les systèmes
d'information géographique (SIG) peuvent faciliter cette représentation
et l'identification des zones insuffisamment desservies. De même, la
modélisation hydraulique, notamment avec des outils tels qu'EPANET,
permet d'étudier les pressions, les débits et différents scénarios de
fonctionnement avant la réalisation des infrastructures
\cite{ref40}, \cite{ref34}.

Dans le cadre de cette étude, le déploiement est considéré comme une
dimension complémentaire de la gestion du réseau de la commune de
Kadutu, notamment pour la représentation spatiale et la simulation de
scénarios d'évolution du réseau.

\section{Surveillance d'un réseau de distribution d'eau potable}
\label{surveillance-dun-reseau-de-distribution-deau-potable}

La surveillance consiste à suivre l'état des infrastructures et
l'évolution des principaux paramètres hydrauliques du réseau. Elle vise
notamment à détecter les fuites, les ruptures, les baisses de pression
et les dysfonctionnements des équipements \cite{ref11}.

Les méthodes traditionnelles reposent principalement sur les inspections
de terrain, les relevés manuels et les signalements des usagers. Elles
peuvent toutefois présenter des limites lorsqu'une anomalie survient
entre deux inspections ou lorsqu'une fuite reste difficilement visible
\cite{ref45}.

L'utilisation de capteurs, de SIG et de plateformes numériques permet
d'améliorer la collecte, la visualisation et l'analyse des informations
relatives au réseau. L'objectif est progressivement de passer d'une
surveillance essentiellement réactive à une surveillance davantage
proactive et fondée sur les données \cite{ref20}, \cite{ref47}.

Dans le présent travail, la surveillance constitue l'une des fonctions
principales du jumeau numérique du réseau de fourniture d'eau potable de
la commune de Kadutu.

\section{Maintenance d'un réseau de distribution d'eau potable}
\label{maintenance-dun-reseau-de-distribution-deau-potable}

La maintenance regroupe les actions visant à maintenir ou à rétablir les
infrastructures dans un état permettant d'assurer la continuité et la
fiabilité du service. Elle peut être corrective, préventive ou
prédictive \cite{ref42}.

La maintenance corrective intervient après une panne, tandis que la
maintenance préventive repose sur des opérations planifiées. La
maintenance prédictive cherche, quant à elle, à exploiter les données
disponibles afin d'identifier les signes annonciateurs d'une défaillance
et de faciliter la planification des interventions \cite{ref45},
\cite{ref49}, \cite{ref47}.

Les fuites constituent un problème important des réseaux d'eau potable.
La surveillance de la pression et du débit, associée à l'analyse des
données, peut contribuer à détecter plus rapidement certaines situations
anormales \cite{ref50}, \cite{ref20}. Dans cette étude, la maintenance
est principalement envisagée sous l'angle de l'aide à la détection des
anomalies et de la priorisation des interventions.

\section{Technologies numériques appliquées à la gestion des réseaux d'eau potable}
\label{technologies-numeriques-appliquees-a-la-gestion-des-reseaux-deau-potable}

La transformation numérique des réseaux d'eau repose sur plusieurs
technologies complémentaires. Les SIG permettent de représenter
spatialement les infrastructures, tandis que l'Internet des objets (IoT)
et les capteurs facilitent la collecte des données de fonctionnement
\cite{ref52}, \cite{ref20}.

Les systèmes de supervision, les bases de données et les solutions
informatiques permettent ensuite de centraliser et d'exploiter les
informations collectées. Les systèmes SCADA sont principalement orientés
vers la supervision et le contrôle, alors que le jumeau numérique ajoute
une dimension de représentation virtuelle, de simulation et d'analyse
\cite{ref54}.

L'intelligence artificielle complète cette chaîne technologique en
permettant l'analyse des données, la détection des anomalies et la
production de prédictions. L'intérêt de ces technologies réside surtout
dans leur intégration au sein d'une architecture cohérente
\cite{ref55}, \cite{ref56}.

\section{Le concept de jumeau numérique (Digital Twin)}
\label{le-concept-de-jumeau-numerique-digital-twin}

Un jumeau numérique est une représentation numérique dynamique d'un
système physique, maintenue en relation avec celui-ci par l'échange de
données. Contrairement à une simple maquette ou à un modèle statique, il
est capable d'évoluer en fonction de l'état observé du système réel
\cite{ref8}, \cite{ref58}.

Le concept est notamment associé aux travaux de Michael Grieves dans le
domaine de la gestion du cycle de vie des produits. Son évolution a été
favorisée par le développement des capteurs, de l'IoT, des bases de
données, de l'informatique en nuage et de l'intelligence artificielle
\cite{ref59}.

Un jumeau numérique repose généralement sur trois éléments essentiels :
le système physique, sa représentation virtuelle et le mécanisme
d'échange de données. Cette organisation permet d'observer l'état du
système, d'analyser son fonctionnement et de simuler différents
scénarios \cite{ref60}.

Dans le domaine de l'eau potable, un jumeau numérique peut représenter
les conduites, les réservoirs, les pompes, les vannes et les autres
composants du réseau. Il peut également intégrer des données de pression,
de débit et de niveau afin de suivre l'état du système et de faciliter
l'analyse des anomalies \cite{ref20}.

Dans le cadre de cette étude, le jumeau numérique est défini comme une
\textbf{représentation virtuelle dynamique du réseau de fourniture d'eau
	potable de la commune de Kadutu}. Il a pour rôle de faciliter la
visualisation du réseau, le suivi de son état, la simulation de certaines
situations et l'aide à la détection des anomalies.

\section{Architecture et fonctionnement d'un jumeau numérique}
\label{architecture-et-fonctionnement-dun-jumeau-numerique}

L'architecture d'un jumeau numérique assure la circulation des données
entre le système physique et sa représentation virtuelle. Elle comprend
généralement une couche physique, des dispositifs de collecte, une
couche de communication, une base de données, un modèle numérique, une
couche d'analyse et une interface de visualisation \cite{ref60}.

Dans un réseau d'eau potable, les capteurs peuvent fournir des mesures
de pression, de débit ou de niveau. Ces données sont transmises,
stockées puis analysées avant d'être présentées sur une interface
permettant aux gestionnaires de suivre l'état du réseau. L'intelligence
artificielle peut être intégrée à cette architecture afin de renforcer
la détection des anomalies et l'aide à la maintenance \cite{ref61}.

Dans cette étude, cette architecture est adaptée au réseau de fourniture
d'eau potable de la commune de Kadutu et sert de base au prototype
développé dans le chapitre III \cite{ref65}.

\section{Intelligence artificielle appliquée à la gestion des réseaux d'eau potable}
\label{intelligence-artificielle-appliquee-a-la-gestion-des-reseaux-deau-potable}

L'intelligence artificielle regroupe des méthodes permettant à un système
informatique d'exploiter des données afin d'effectuer des tâches
d'analyse, de classification ou de prédiction. Parmi ses principales
branches figure l'apprentissage automatique (\emph{Machine Learning}),
qui permet à un modèle d'apprendre des relations à partir de données
\cite{ref67}, \cite{ref69}.

L'apprentissage automatique comprend notamment l'apprentissage supervisé,
non supervisé et par renforcement. Pour la gestion des réseaux d'eau,
plusieurs algorithmes peuvent être utilisés, tels que les arbres de
décision, les forêts aléatoires, les SVM et les réseaux de neurones,
selon la nature des données et le problème étudié \cite{ref70},
\cite{ref71}.

L'intelligence artificielle peut notamment être utilisée pour détecter
des anomalies, identifier des comportements inhabituels, analyser les
variations de pression et de débit et contribuer à l'anticipation de
certaines défaillances \cite{ref73}.

Dans le présent travail, l'IA est intégrée au jumeau numérique pour
analyser les données disponibles et contribuer à la détection de
certaines anomalies du réseau de la commune de Kadutu. Son utilisation
reste un outil d'aide à la décision et ne se substitue pas à
l'expertise des gestionnaires.

\section{Intégration du jumeau numérique et de l'intelligence artificielle}
\label{integration-du-jumeau-numerique-et-de-lintelligence-artificielle}

Le jumeau numérique fournit une représentation dynamique du système
physique, tandis que l'intelligence artificielle apporte des capacités
d'analyse et de détection. Leur intégration permet donc de transformer
une simple représentation numérique en un système capable d'exploiter les
données pour améliorer la surveillance et la maintenance
\cite{ref60}, \cite{ref61}.

Dans un réseau d'eau potable, les données relatives à la pression, au
débit ou à d'autres paramètres peuvent être analysées afin d'identifier
des comportements inhabituels. Les résultats peuvent ensuite être
présentés sous forme d'alertes ou d'informations utiles à la
priorisation des interventions.

Cette intégration présente toutefois des limites, notamment la
disponibilité des données, leur qualité, l'interopérabilité des
technologies et les contraintes liées aux ressources techniques et
financières \cite{ref60}. Ces contraintes sont particulièrement
importantes dans les contextes où les infrastructures numériques restent
limitées.

Dans le cadre de cette recherche, l'intégration du jumeau numérique et de
l'intelligence artificielle vise principalement à améliorer la
surveillance, la détection des anomalies et l'aide à la maintenance du
réseau de fourniture d'eau potable de la commune de Kadutu.

\section{Revue des travaux antérieurs}
\label{revue-des-travaux-anterieurs}

Les travaux antérieurs permettent d'identifier les principales
approches déjà utilisées dans le domaine des jumeaux numériques, de
l'intelligence artificielle et de la gestion intelligente des réseaux
d'eau potable. Les recherches montrent que le jumeau numérique peut
améliorer la représentation, la surveillance et l'analyse des
infrastructures, tandis que l'intelligence artificielle contribue à la
détection des anomalies et à l'analyse prédictive \cite{ref8},
\cite{ref17}.

Tao et al. ont notamment proposé une architecture générale du jumeau
numérique fondée sur l'intégration des objets connectés, des données et
des modèles numériques. D'autres travaux ont montré l'intérêt du
Machine Learning pour l'analyse des réseaux d'eau et le développement
des systèmes de gestion intelligente \cite{ref17}, \cite{ref76}.

En Afrique subsaharienne, plusieurs initiatives ont exploré l'utilisation
des capteurs, des SIG et des plateformes numériques pour améliorer la
surveillance des réseaux. Toutefois, les solutions intégrant
simultanément jumeau numérique, intelligence artificielle et IoT restent
encore limitées \cite{ref77}.

En République Démocratique du Congo, les travaux portant spécifiquement
sur cette combinaison technologique sont encore peu nombreux. Les
recherches existantes concernent principalement l'accès à l'eau, les
performances des infrastructures et les pertes d'eau. Cette situation
laisse apparaître un espace de recherche pour une solution adaptée au
contexte de la commune de Kadutu \cite{ref78}, \cite{ref79}.

La présente étude se positionne donc dans cette continuité en proposant
un \textbf{jumeau numérique intégrant des techniques d'intelligence
	artificielle} spécifiquement adapté au réseau de fourniture d'eau
potable de la commune de Kadutu. Le système vise principalement à
faciliter la représentation du réseau, la surveillance, la détection des
anomalies et l'aide à la maintenance.

\paragraph{Tableau 1.1 : Synthèse comparative des principaux travaux antérieurs}

\begin{longtable}[]{@{}
		>{\raggedright\arraybackslash}p{(\columnwidth - 10\tabcolsep) * \real{0.1483}}
		>{\raggedright\arraybackslash}p{(\columnwidth - 10\tabcolsep) * \real{0.0792}}
		>{\raggedright\arraybackslash}p{(\columnwidth - 10\tabcolsep) * \real{0.1387}}
		>{\raggedright\arraybackslash}p{(\columnwidth - 10\tabcolsep) * \real{0.1982}}
		>{\raggedright\arraybackslash}p{(\columnwidth - 10\tabcolsep) * \real{0.1950}}
		>{\raggedright\arraybackslash}p{(\columnwidth - 10\tabcolsep) * \real{0.2407}}@{}}
	\toprule
	Auteur(s) & Année & Domaine d'étude & Contribution principale &
	Limites & Apport de la présente étude \\
	\midrule
	\endhead
	
	Grieves & 2014 & Jumeau numérique &
	Formalisation du concept de Digital Twin &
	Approche générale &
	Application au réseau d'eau potable de Kadutu \\
	
	Tao et al. & 2019 & Industrie 4.0 &
	Architecture générale du jumeau numérique &
	Orientation industrielle &
	Adaptation aux infrastructures hydrauliques \\
	
	Fuller et al. & 2020 & Digital Twin &
	Technologies, défis et perspectives &
	Peu d'applications aux pays en développement &
	Prise en compte du contexte de Kadutu \\
	
	Abu-Mahfouz et al. & 2020 & Smart Water Grid &
	Gestion intelligente des réseaux d'eau &
	Intégration limitée du jumeau numérique &
	Association du jumeau numérique et de l'IA \\
	
	Rasheed et al. & 2020 & Digital Twin &
	Analyse des modèles prédictifs &
	Peu orienté vers les réseaux d'eau &
	Accent sur la détection des anomalies et la maintenance \\
	\bottomrule
\end{longtable}